\chapter{The cost model: what a desert costs, and how predictably}
\label{ch:costmodel}

\begin{chapterabstract}
A partial-cover construction is a lottery: each progression element
survives only if every residual complement happens to be composite. This
chapter turns that lottery into a measured science. The success density
is $e^{-\Ee}$ with $\Ee=|U|\cdot\mathrm{boost}/\ln\N$; we derive the
model, refine it with the exact survivor-count distribution, and
validate it to $\pm0.1$ nats across five scales and $1.2\cdot10^9$
scanned candidates --- the dice are exactly as loaded as predicted, every
time. On top of the model we build the economics: the digit-budget
identity that splits $\ln\N$ into modulus nats and search nats, the
digit--compute exchange rate, the two binding regimes of the bounded
games, the boost-versus-kills equilibrium that sets cover size, and the
frontier curves $\Ee(D)$ and $\Ee(Q)$ that price every record before it
is attempted. Two consequences frame the rest of the thesis: a
million-desert below $200$ digits exists in abundance heuristically yet
costs $10^{85}$ primality tests to exhibit by covers --- a
construction-versus-existence gap of fifty-nine orders of magnitude
beyond cosmological --- and the calibrated truth for $T(100\,000)$ sits
near $10^{53}$, a floor our $150$-digit record approaches only in
logarithm.
\end{chapterabstract}

\section{The lottery}\label{sec:lottery}

Fix a partial cover: odd prime moduli $\mathcal R$ with residues
$a_r$, parity, the modulus $M=2\prod_{r\in\mathcal R}r$, and the
residual set $U$ of uncovered odd primes below the target $Q$. Every
candidate is $\N=\N_0+kM$, and $k$ wins exactly when all $|U|$
residual complements $\N-q$ ($q\in U$) are composite
(Proposition~\ref{prop:partial}). Losing tickets are cheap to recognise:
scan the residual complements and stop at the first probable prime.
Winning tickets are certified as in Chapter~\ref{ch:intro}.

Two structural facts make the lottery analysable. First, for residual
$q$ the complement $\N-q$ is automatically coprime to $2$ and to every
cover modulus: $q\not\equiv a_r=\N\bmod r$ means $r\nmid\N-q$. The cover
therefore pre-sieves its own residuals --- the roughness bonus of
Remark~\ref{rem:rough} in its useful form. Second, the residual
complements at a given $k$ are distinct integers of size $\approx\N$
whose only shared structure is exactly that coprimality (a fact made
precise, and load-bearing, in Chapter~\ref{ch:closure}).

\section{The density model}\label{sec:density}

A random integer near $\N$ is prime with probability about $1/\ln\N$.
Conditioning on coprimality to a set $S$ of primes divides the candidate
space by $\prod_{p\in S}(1-1/p)$, so the conditional primality
probability for one residual complement is
\[
p_1\;=\;\frac{\mathrm{boost}}{\ln\N},
\qquad
\mathrm{boost}\;=\;2\prod_{r\in\mathcal R}\frac{r}{r-1},
\]
the factor $2$ from oddness. When the cover contains every prime up to
$r_{\max}$, Mertens' theorem gives
$\mathrm{boost}\approx2e^{\gamma}\ln r_{\max}$; real covers skip some
primes and the product is computed exactly. Treating the $|U|$ residual
primality events as independent (the whole content of that assumption is
examined in Chapter~\ref{ch:closure}; empirically it holds to
$\le0.1$ nats), the probability that a candidate wins is
\[
\prod_{q\in U}(1-p_1)\;\approx\;e^{-|U|\,p_1}\;=\;e^{-\Ee},
\qquad
\boxed{\;\Ee\;=\;\frac{|U|\cdot\mathrm{boost}}{\ln\N}\;}
\]
and the expected number of progression elements to the first hit is
$e^{\Ee}$. The exponent $\Ee$ is the single number that decided every
campaign in this thesis: it is the negative log of the per-candidate
success density, measured in nats.

\paragraph{The exact refinement.} In practice the search engine sieves
each block of $k$ by all primes up to a depth $B_{\mathrm{sieve}}$
(beyond the cover moduli), so at each $k$ only $a_k\le|U|$ complements
remain \emph{alive}, and a surviving complement is prime with the deeper
conditional probability $\hat p$. The exact per-$k$ density is
$(1-\hat p)^{a_k}$, and the honest exponent is
\[
\Ee_{\mathrm{exact}}\;=\;-\ln\;\mathbb E_k\bigl[(1-\hat p)^{a_k}\bigr],
\]
the expectation over the sieve's survivor-count distribution (measured
mean $\approx316$, standard deviation $\approx14$ for a typical
$|U|=750$ cover at $199$ digits, with $\hat p\approx0.059$). Because the
average of an exponential exceeds the exponential of the average, the
fluctuation in $a_k$ is worth a small bonus, $e^{\hat
p^2\sigma_a^2/2}\approx0.2$--$0.3$ nats; the naive formula undercounts
by exactly that much, and the measured densities below confirm the
exact form. The same survivor counts drive the search-order
optimizations of Chapter~\ref{ch:records}.

\section{Validation: five scales, \texorpdfstring{$1.2\cdot10^9$}{1.2e9} candidates}\label{sec:validation}

The model was validated before it was trusted, at every scale it was
later used, by comparing the predicted exponent against
$\hat\Ee=-\ln(\text{measured survival rate})$ on fully logged scans.

\begin{table}[h]
\centering\small
\caption{Density-model validation. ``Predicted'' is
$\Ee$ (exact form where computed); ``measured'' is $\hat\Ee$ from the
logged scan. The two toy deserts are small constructions built to be
exhaustively scannable.}\label{tab:validation}
\begin{tabular}{lrrr}
\toprule
construction & digits of $\N$ & predicted $\Ee$ & measured $\hat\Ee$\\
\midrule
toy desert & 26 & 12.53 & 12.51\\
toy desert & 62 & 13.52 & 13.57\\
dual-game record cover & 197 & 18.20 & 18.2\\
threshold record cover & 195 & 17.51 & 17.4\\
GPU threshold cover & 150 & 24.91 & (one hit at $\lambda=0.63$)\\
mega-desert cover & 2\,692 & 6.60 & 6.56\\
\bottomrule
\end{tabular}
\end{table}

Across three fully instrumented campaign covers the running discrepancy
$|\hat\Ee-\Ee|$ stayed below $0.1$ nats over $1.2\cdot10^9$ scanned
progression elements, with per-slice survivor rates within $\pm2\%$ of
prediction across more than $500$ logged slices. The $150$-digit row
deserves its asterisk explained: that campaign scanned its entire
$k$-round $[0,7.5\cdot10^{10})$ end to end ($\sim2.1\cdot10^{12}$
Fermat tests) and collected exactly one hit against a Poisson mean of
$0.63$ --- an unremarkable draw, and the model stayed unfalsified at the
largest scale ever pushed. This validation is what licenses every
forecast in the thesis: when we say a rung costs $e^{27.3}$ candidates,
that number has survived contact with $10^9$-scale reality repeatedly.

We flag the honest converse: the validation also \emph{bounds} every
dream of beating the lottery. Any exploitable correlation between
residual complements, any hidden compositeness bias in constructed $\N$,
would show up as $\hat\Ee<\Ee$; over $1.2\cdot10^9$ candidates nothing
beyond $0.1$ nats ever did. Chapter~\ref{ch:closure} explains why not.

\section{The digit-budget identity and the exchange rate}\label{sec:budget}

Take logarithms of $\N\approx kM$:
\[
\underbrace{\ln\N}_{\text{record size (nats)}}
\;=\;
\underbrace{\ln M}_{\text{modulus nats}}
\;+\;
\underbrace{\ln k}_{\text{search nats}}.
\]
Modulus nats are spent at design time on congruences (each modulus $r$
costs $\ln r$ and kills a residue class of offsets); search nats are
spent at run time scanning the progression ($\ln k^*\approx\Ee$ at the
expected hit). The two currencies buy the same digits, and the whole
craft of the bounded games is arbitraging them:

\begin{itemize}
\item \textbf{Search nats cost only compute.} A 4-core CPU pipeline
demonstrated $\ln k\approx20$; the GPU fleet of
Chapter~\ref{ch:records} afforded $\ln k\approx25$. Throughput converts
to record digits at the model's slope --- which is why the single
largest threshold-game jump in this thesis ($179\to150$ digits) was
purely an engineering event: the CUDA engine's $430\times$ throughput is
$\ln430\approx6.1$ nats, precisely the exponent gap ($\Ee=24.9$ vs
$\approx19$) between the two rungs.
\item \textbf{The exchange rate is the frontier slope.} Near $200$
digits, $d\Ee/dD\approx-0.10$ to $-0.19$ per digit (so $10\times$
compute buys $12$--$23$ digits there); by $150$ digits the slope is
$\approx-0.32$, and near $100$ digits $\approx-0.5$ ($10\times$ buys
only $\sim5$). The certified ladder measured it directly: each
$\sim7$-digit rung of the $199\to179$ descent cost $\approx2.5\times$
the previous search.
\item \textbf{Cover quality is the third currency.} At fixed $\ln M$,
better residue choices shrink $|U|$; each $1\%$ of $|U|$ is worth
$e^{0.01\cdot\Ee}$ of search. Chapter~\ref{ch:closure} shows the greedy
covers are locally perfect and globally within an integrality gap worth
at most a few nats --- the least solved, most valuable margin in the
game.
\end{itemize}

\section{Two binding regimes, and the equilibrium that sizes covers}\label{sec:regimes}

\paragraph{Compute-bound versus range-bound.} With no size cap, the only
constraint is affordable search: pick the cover minimising
$\ln M+\Ee$ and scan. Under a cap $\N<B$ (the bounded games), the
progression only holds $k_{\max}=B/M$ candidates, and the binding
constraint can flip from compute to \emph{range}. In the range-bound
regime the optimum counterintuitively \emph{shrinks} the modulus below
its unconstrained value: each nat removed from $\ln M$ buys a full nat
of $k$-range while costing only $\approx0.2$ nats of exponent (the
frontier slope in disguise). This was derived, then confirmed twice in
campaign post-mortems: raw GPU throughput pointed at a $200$-digit
budget target without re-optimising $M$ failed to beat a tuned CPU
cover, and the tuned cover's own success confirmed the shrunken-modulus
optimum.

\paragraph{The boost-versus-kills equilibrium.} What stops a cover from
growing forever? Adding a modulus $r$ with its best residue class kills
some number of residuals ($\mathrm{kills}\approx|U|/(r-1)$ on random
structure, less at the picked-over endgame), lowering $\Ee$ by
$\mathrm{kills}\cdot\mathrm{boost}/\ln\N$; but it also multiplies the
boost by $r/(r-1)$, \emph{raising} every surviving residual's primality
odds and hence $\Ee$ by $|U|\cdot\mathrm{boost}/((r-1)\ln\N)$. The two
cancel at
\[
\mathrm{kills}\;\approx\;\frac{|U|}{r-1},
\]
and past that point new moduli are net-zero or worse. Measured at
$B=10^{200}$, $Q=122\,000$: adding $r=461$ to the standing $88$-modulus
cover kills $3.5$ residuals ($-0.084$ nats) and boosts the survivors
($+0.082$ nats) --- net $+0.00$. Two consequences. First, at large
budgets the cover is \emph{not} budget-limited ($\ln M$ used $437$ of
$\sim458$ available nats in that instance); it sits at the equilibrium,
which is why leftover digit budget buys nothing and why the frontier is
flat against re-derivation --- independent greedy runs at the same $Q$
land within $10^{-3}$ nats of each other. Second, the binding constraint
crosses over: at $Q=122\,000$ under $199$ digits the \emph{budget}
binds, while at $Q=200\,000$ the \emph{equilibrium} binds. Both regimes
were mapped before the campaigns that used them.

\section{The frontier curves}\label{sec:frontiers}

Everything above compresses into two curves, computed by building actual
covers (lazy-greedy, digit-capped, annealed) and evaluating $\Ee$
exactly.

\begin{table}[h]
\centering\small
\caption{The digit frontier $\Ee(D)$ at fixed threshold $Q=100\,003$
(greedy digit-capped covers). ``Candidates'' is the expected search
$e^{\Ee}$.}\label{tab:ED}
\begin{tabular}{rrrrr}
\toprule
$D$ (digits) & moduli & $|U|$ & $\Ee$ & candidates\\
\midrule
100 & 52 & 973 & 41.3 & $10^{17.9}$\\
120 & 60 & 899 & 32.8 & $10^{14.2}$\\
130 & 64 & 862 & 29.5 & $10^{12.8}$\\
140 & 68 & 830 & 26.8 & $10^{11.6}$\\
150 & 72 & 802 & 24.4 & $10^{10.6}$\\
160 & 76 & 773 & 22.3 & $10^{9.7}$\\
170 & 79 & 754 & 20.6 & $10^{8.9}$\\
195 & 89 & 692 & 16.8 & $10^{7.3}$\\
\bottomrule
\end{tabular}
\end{table}

\begin{table}[h]
\centering\small
\caption{The threshold frontier $\Ee(Q)$ at fixed budget $199$ digits
(same $88$-modulus digit budget, greedy + coordinate ascent + kicks;
boost $11.01$). The last column prices each rung against its
predecessor.}\label{tab:EQ}
\begin{tabular}{rrrrr}
\toprule
$Q$ & $|U|$ & $\Ee$ & candidates & $\times$ previous\\
\midrule
122\,000 & 867 & 20.83 & $1.1\cdot10^{9}$ & ---\\
135\,000 & 1\,004 & 24.12 & $3.0\cdot10^{10}$ & $\times27$\\
150\,000 & 1\,138 & 27.34 & $7.5\cdot10^{11}$ & $\times25$\\
175\,000 & 1\,335 & 32.07 & $8.5\cdot10^{13}$ & $\times113$\\
200\,000 & 1\,569 & 37.69 & $2.4\cdot10^{16}$ & $\times280$\\
\bottomrule
\end{tabular}
\end{table}

Two planning corollaries. Because rung spacing is geometric, climbing
\emph{every} rung of a ladder costs only $\approx4\%$ more than jumping
straight to the top --- and banks a certified record, plus a fresh
calibration of $\hat\Ee$, at each step. And each curve exposes its
regime: Table~\ref{tab:ED} is the compute-bound threshold game (digits
fall as affordable $\Ee$ rises), Table~\ref{tab:EQ} the range-bound
budget game (the target rises until the equilibrium chokes it).

\section{Construction versus existence: the million-desert study}\label{sec:million}

How far can the paradigm be pushed? As a stress test we priced the
\emph{million desert}: an even $\N<10^{200}$ with $\g(\N)>10^6$, i.e.\
all $78\,498$ primes $q\le10^6$ blocked at once. Note the conjectural
ceiling \eqref{eq:GLR} evaluates to
$(\ln\N)^2\ln\ln\N\approx1.3\cdot10^6$ exactly at $200$ digits: this
target sits \emph{at} the extremal constant of the
Granville--van de Lune--te Riele law, not beyond it. And existence is
not the issue: for unstructured even $\N$ near $10^{200}$ the
all-composite exponent is $\Ee_{\mathrm{random}}\approx445$, so the
expected count of qualifying integers below $10^{200}$ is
$10^{200}e^{-445}\approx10^{6.5}$ --- they exist in abundance,
heuristically.

Exhibiting one is another matter. The frontier $\Ee(D)$ at $Q=10^6$:

\begin{table}[h]
\centering\small
\caption{The million-desert frontier (lazy-greedy digit-capped covers,
$Q=10^6$).}\label{tab:million}
\begin{tabular}{rrrrr}
\toprule
$D$ (digits) & moduli & $|U|$ & $\Ee$ & search cost $e^{\Ee}$\\
\midrule
199 & 91 & 8\,157 & 196.7 & $10^{85.4}$\\
300 & 127 & 7\,311 & 124.3 & $10^{54.0}$\\
400 & 161 & 6\,728 & 89.3 & $10^{38.8}$\\
600 & 226 & 5\,865 & 54.9 & $10^{23.8}$\\
800 & 288 & 5\,241 & 38.2 & $10^{16.6}$\\
1000 & 347 & 4\,741 & 28.4 & $10^{12.3}$\\
1300 & 435 & 4\,100 & 19.6 & $10^{8.5}$\\
1600 & 519 & 3\,588 & 14.2 & $10^{6.2}$\\
2400 & 736 & 2\,492 & 6.9 & $10^{3.0}$\\
\bottomrule
\end{tabular}
\end{table}

The $2400$-digit row reproduces the regime of the certified mega-desert
of Chapter~\ref{ch:gaps} ($\Ee\ 6.9$ vs its measured $6.56$) --- the
ladder is consistent end to end. Reading off feasibility: a
record-scale search ($1.4\cdot10^8$ candidates) needs
$\sim$$1\,470$ digits; a heavy single box ($10^9$) $\sim$$1\,250$; a
distributed fleet ($7\cdot10^{10}$) $\sim$$1\,120$; and a cosmological
budget ($10^{12}$ tests per second for the age of the universe,
$\Ee\le68$) still needs $\sim$$520$ digits. A verified $1\,250$-digit
blueprint ($420$ moduli, $|U|=4\,205$, $\Ee=20.74$, about $235$
core-days) is committed in the artifact repository as the realistic
version of this target.

At $199$ digits the wall is structural. The best $199$-digit system
found --- $91$ moduli, residue-annealed --- leaves $|U|=8\,146$ and
$\Ee=196.4$; annealing recovers only $\sim0.3\%$ of $\Ee$. Structural in
the sieve sense, not the counting sense: the $91$ best classes hold
$148\,290$ slots for $78\,497$ targets, a $1.89\times$ supply, so a
union bound alone forbids nothing --- but prime residue classes overlap
like independent sieves, coverage stalls near $90\%$, and $60\,000$
anneal moves reclaim $11$ residuals out of $8\,157$. Joint optimization
of representative and coverage (Chapter~\ref{ch:limits}) moves
$\Ee$ from $196$ to $\approx194$. So: about $10^{6.5}$ million-deserts
exist below $10^{200}$, and the cheapest known path to \emph{point at
one} costs $10^{85}$ primality tests --- a construction-versus-existence
gap of fifty-nine orders of magnitude beyond a cosmological compute
budget. That gap, in miniature, is also the story of the sub-$100$-digit
threshold game (Chapter~\ref{ch:limits}); the million-desert study is
the thesis's cleanest exhibit that \emph{knowing something exists and
being able to exhibit it are different sciences}.

\section{Poisson discipline: negative results without excuses}\label{sec:poisson}

The lottery is memoryless: conditioned on no hit so far, the expected
remaining wait is unchanged. That single fact, taken seriously, is most
of campaign discipline.

\begin{itemize}
\item \textbf{Spread.} Over the $n=7$ fully-logged searches that ended
in hits, first hits landed at $0.25\times$--$3.0\times$ the expected
$e^{\Ee}$ (median $0.5\times$) on a consistent consumed-candidates
basis. The exponential clock is merciless about variance: plan budgets
at $3\times\Ee$, in writing, before starting.
\item \textbf{Misses are data.} Campaign 1 of the PhD track retired
hitless at $3.18\cdot10^8$ elements --- survival probability $0.30$
under the calibrated model, an unlucky but ordinary draw; per-slice
$\hat p$ and survivor counts matched prediction throughout, ruling out
implementation error. Campaign 2 stopped at $5.6\cdot10^8$ elements
($P\approx0.22$). The model was never adjusted after the fact to
absorb a miss; it never needed to be.
\item \textbf{The tail teaches.} The first-generation-beating
$199$-digit dual record (Chapter~\ref{ch:records}) hit at $99.9\%$ of
its committed range, $3.0\times$ the selection-adjusted expectation ---
the outlier that taught the $3\times$ rule.
\item \textbf{Ratchet, don't grand-slam.} Under fixed compute, a ladder
of certified rungs dominates a single deep campaign: geometric rung
spacing makes the total cost $\approx1.04\times$ the top rung
(\S\ref{sec:frontiers}), each rung banks a record and a calibration,
and optimal stopping becomes arithmetic. The thesis learned this the
expensive way --- a breadth-first 400-variant campaign capped its best
possible outcome and exhausted hitless while a rival single-progression
depth-first design (unbounded runway, bad luck costs digits rather than
termination) banked the record. Both post-mortems are in
Chapter~\ref{ch:records}.
\end{itemize}

\section{Calibrating the truth: how far are the records from $T(Q)$ itself?}\label{sec:truth}

The exhaustive data fix the constant in the conjectural
law~\eqref{eq:GLR}: $\g\le9\,781$ for all even $\N\le4\cdot10^{18}$
gives $C\approx1.4$ in
$\g_{\max}(X)\approx C\,\ln^2\!X\,\ln\ln X$. Inverting at
$Q=10^5$ puts the \emph{true} $T(100\,000)$ near $10^{53}$ ---
plus-or-minus a few in the exponent from the unknown constant and
residue-luck tails --- i.e.\ a $53$-digit integer, unreachable by any
imaginable exhaustive verification (it would take $\sim10^{52}$ scans)
and, Chapter~\ref{ch:limits} will argue, essentially unreachable by
construction too. Our $150$-digit record is thus about $2.9\times$ the
truth \emph{in logarithm}. For the budget game the truth is closer:
$R(10^{200})\approx C\cdot(460.5)^2\ln(460.5)\approx1.8$ million,
against a certified $119\,419$ --- and there the gap is
compute-and-covers, not existence.

For threshold games the model collapses to a useful pair of envelope
identities (cover-quality factor $\varphi\approx0.60$ for current
optimizers; $\ln\mathcal C$ the affordable search exponent):
\begin{align*}
\mathrm{digits}(\N)\;&\gtrsim\;\frac{0.57\,\varphi\,\pi(Q)}{\ln\mathcal C}
&&\text{(compute-bound wall)},\\
\mathrm{digits}(\N)\;&\gtrsim\;\frac{\sqrt{c\,\pi(Q)}}{2.3},\ c\in[0.8,1.3]
&&\text{(absolute floor)},
\end{align*}
the floor being the regime where the ``construction'' has degenerated
into nature's own ensemble scan --- which is why the floor
($38$--$49$ digits at $Q=10^5$) and the truth ($\sim53$ digits) nearly
coincide. Everything between the frontier tables and that floor is the
subject of Chapters~\ref{ch:limits} and~\ref{ch:closure}: how much is
buyable with compute, how much with better covers, and how much is
provably out of reach.
